% Owner: TV1 (P5.1a) -- Reviewer: TV2
\section{Related Work}
\label{sec:related_work}

\textbf{Minimax and alpha-beta pruning.} Shannon~\cite{shannon1950programming} framed computer game playing as a depth-limited Minimax search over the game tree, with a static evaluation function applied at the horizon. Alpha-beta pruning returns the same value as Minimax while skipping branches that cannot change the result. Knuth and Moore~\cite{knuth1975analysis} proved its correctness and showed that, with perfect move ordering, a uniform tree of branching factor $b$ and depth $d$ needs only about $b^{\lceil d/2\rceil}+b^{\lfloor d/2\rfloor}$ leaf evaluations instead of $b^d$; with random ordering the savings are much smaller. The quality of move ordering is therefore the main lever on the cost of alpha-beta, which is what Table~T1 measures on $3\times 3$ Tic-Tac-Toe.

\textbf{Solving games.} Small games can be solved exactly: $3\times 3$ Tic-Tac-Toe is a draw under perfect play, which makes it a convenient ground truth for testing whether a sampling-based player ever loses. Allis~\cite{allis1994searching} studied which games can be solved by search and, using threat-space and proof-number search, showed that free-style Gomoku on the $15\times 15$ board is a win for the first player. Such results rely on game-specific threat analysis; general players on large boards must still trade search depth against evaluation quality.

\textbf{Bandits and UCT.} The multi-armed bandit problem formalizes the trade-off between exploiting the best-looking action and exploring the others. Lai and Robbins~\cite{lai1985asymptotically} showed that regret must grow at least logarithmically, and Auer et al.~\cite{auer2002finite} gave UCB1, a simple index policy that achieves logarithmic regret uniformly over time. Monte Carlo tree search grew out of Monte Carlo evaluation in Go: Coulom~\cite{coulom2006efficient} combined random simulations with an incrementally built tree, and Kocsis and Szepesv\'ari~\cite{kocsis2006bandit} proposed UCT, which applies UCB1 at every node of the tree and converges to the Minimax value. Coquelin and Munos~\cite{coquelin2007bandit} showed that this convergence can be very slow on unfavorable trees. Browne et al.~\cite{browne2012survey} survey the many MCTS variants, including alternative selection formulas, final-move selection rules, and domain-specific rollout policies, which are the design choices studied in Sections~\ref{sec:methodology} and~\ref{sec:results}.

\textbf{Search with learned knowledge.} AlphaGo~\cite{silver2016mastering} replaced hand-written rollout and evaluation knowledge with deep policy and value networks trained from expert games and self-play, and defeated a professional Go player. AlphaGo Zero~\cite{silver2017mastering} removed human data and rollouts entirely, guiding MCTS only by a network trained through self-play. Our rollout policies L0--L4 and the pattern-based evaluation used for alpha-beta are hand-crafted counterparts of these policy and value functions, which we return to in Section~\ref{sec:discussion}.
